\documentclass[14pt, a4paper]{extarticle} % или article, если не нужен увеличенный базовый шрифт
\usepackage[utf8]{inputenc}
\usepackage[T2A]{fontenc}
\usepackage[russian]{babel}
\usepackage{array}
\usepackage{graphicx}
\usepackage{ulem}
\usepackage{geometry}
\geometry{top=2cm, bottom=2cm, left=2cm, right=2cm}
\usepackage{tabularx}
\usepackage{amsmath}
\usepackage{amssymb}
\usepackage{multirow}
\usepackage{booktabs}
\usepackage{float}
\usepackage{tabularx}
\usepackage{booktabs}
\usepackage{siunitx}
\usepackage{indentfirst}
\usepackage{placeins}
\usepackage{lipsum}
\usepackage{longtable}
\usepackage{csquotes}
\usepackage[style=gost-numeric,sorting=none]{biblatex}
\usepackage{enumitem}
\usepackage{caption}

\captionsetup[table]{position=top, singlelinecheck=false, justification=raggedleft}



\addbibresource{references.bib}
\usepackage{titlesec}
\titleformat{\chapter}[display]
    {\normalfont\huge\bfseries}
    {\chaptertitlename\ \thechapter}
    {20pt}
    {\Huge}
\usepackage{booktabs}
\usepackage{makecell}
\graphicspath{{Image/}}

\geometry{
    a4paper,
    total={170mm,257mm},
    left=20mm,
    top=20mm,
}

\newcommand{\specialcell}[2][c]{%
    \begin{tabular}[#1]{@{}c@{}}#2\end{tabular}%
}

\usepackage{etoolbox}
\patchcmd{\thebibliography}
  {\chapter*{\bibname}}
  {\section*{\centering СПИСОК ЛИТЕРАТУРЫ}}
  {}{}
\patchcmd{\thebibliography}
  {\addcontentsline{toc}{chapter}{\bibname}}
  {\addcontentsline{toc}{section}{СПИСОК ЛИТЕРАТУРЫ}}
  {}{}


\usepackage{setspace}
\onehalfspacing


% Для подчеркивания текста, аналогичного .underline в Word
\newcommand{\und}[1]{\uline{#1}}

\begin{document}

% Титульный лист
\begin{titlepage}
    \centering
    \noindent
    \begin{tabular}{@{}p{0.25\textwidth} p{0.7\textwidth}@{}}
        \raisebox{-\totalheight}{\includegraphics[width=0.8in, height=0.9in]{imgs/BMSTU.jpg}} &
        \textbf{Министерство науки и высшего образования Российской Федерации} \\
        & \\
        & \textbf{Федеральное государственное бюджетное образовательное} \\
        & \textbf{учреждение} \\
        & \\
        & \textbf{высшего образования} \\
        & \\
        & \textbf{«Московский государственный технический университет} \\
        & \\
        & \textbf{имени Н.Э. Баумана} \\
        & \\
        & \textbf{(национальный исследовательский университет)»} \\
        & \\
        & \textbf{(МГТУ им. Н.Э. Баумана)} \\
    \end{tabular}

    \vspace{1cm}
    ФАКУЛЬТЕТ \und{«СПЕЦИАЛЬНОЕ МАШИНОСТРОЕНИЕ»}

    \vspace{0.5cm}
    КАФЕДРА \und{«РАКЕТНЫЕ И ИМПУЛЬСНЫЕ СИСТЕМЫ» (СМ-6)}

    \vspace{2cm}
    \textbf{РАСЧЕТНО-ПОЯСНИТЕЛЬНАЯ ЗАПИСКА}

    \vspace{1cm}
    \textbf{\textit{К КУРСОВОЙ РАБОТЕ}}

    \vspace{1cm}
    \textbf{\textit{НА ТЕМУ:}}

    \vspace{0.5cm}
    \textbf{\textit{\und{«Баллистическое проектирование артиллерийских орудий»}}}

    \vspace{1.5cm}
    \textbf{Вариант 24}

    \vspace{2cm}
    \begin{flushleft}
        Студент \und{СМ6-72} \hfill \underline{\hspace{5cm}} \hfill \und{М.В. Ерофеев} \\
        \hspace{4.8cm} (Группа) \hspace{2.3cm} (Подпись, дата) \hspace{1.8cm} (И.О.Фамилия) \\[1cm]
        Руководитель курсовой работы \hfill \underline{\hspace{5cm}} \hfill \und{В.А. Федулов} \\
        \hspace{7.5cm} (Подпись, дата) \hspace{1.8cm} (И.О.Фамилия)
    \end{flushleft}

    \vfill
    \centering
    \textit{2025 г.}

    \newpage % Начало второй страницы (Задание)

    \centering
    \textbf{Министерство науки и высшего образования Российской Федерации}

    \textbf{Федеральное государственное бюджетное образовательное учреждение}

    \textbf{высшего образования}

    \textbf{«Московский государственный технический университет имени Н.Э. Баумана}

    \textbf{(национальный исследовательский университет)»}

    \textbf{(МГТУ им. Н.Э. Баумана)}

    \vspace{1cm}
    УТВЕРЖДАЮ

    \vspace{0.5cm}
    Заведующий кафедрой \und{СМ6} \\
    \hspace{1.5cm} (Индекс) \\[0.5cm]
    \underline{\hspace{5cm}} \und{В.М.Кашин} \\
    \hspace{0.5cm} (И.О.Фамилия) \\[0.5cm]
    « \underline{\hspace{1cm}} » \underline{\hspace{3cm}} 20 \underline{\hspace{1cm}} г.

    \vspace{2cm}
    \textbf{ЗАДАНИЕ}

    \textbf{на выполнение курсовой работы}

    \vspace{1cm}
    \begin{flushleft}
        по дисциплине \und{Газовая динамика} \hfill \underline{\hspace{7cm}} \\[0.5cm]
        Студент группы \und{СМ6-72} \\[0.5cm]
        \underline{\hspace{12cm}} \und{Ерофеев Максим Викторович} \\
        \hspace{5cm} (Фамилия, имя, отчество) \\[0.5cm]
        Тема курсовой работы \und{Баллистическое проектирование артиллерийских орудий} \\
        \und{на сжатом газе} \\[0.5cm]
        Направленность КР (учебная, исследовательская, практическая, производственная, др.) \\
        \und{Учебная} \hfill \underline{\hspace{7cm}} \\[0.5cm]
        Источник тематики (кафедра, предприятие, НИР) \und{Кафедра} \hfill \underline{\hspace{5cm}} \\[0.5cm]
        График выполнения работы: 25\% к \und{4} нед., 50\% к \und{8} нед., 75\% к \und{12} нед., 100\% к \und{14} нед. \\[1cm]
        \textbf{\textit{Задание}} \und{Найти оптимальные параметры ствола и условий заряжания при решении обратной задачи внутренней баллистики. При этом $d = 85$ мм, $q =$ 5 кг, $ v_{\text{pm}} = 950 $ м/с, тип орудия – НР, тип математической модели – квазиодномерная (КМ), критерий оптимальности – $Z_{\text{B1}}$, $p^{\text{max}}_{\text{m}}$ = 390 МПа, $l^{\text{max}}_{\text{m}}$ = 65 ед.d, $v_{\text{pm-50}}$ = 850 м/с, $p_{\text{mz+5}}0 = $ 180 МПа.} \\[1cm]
        \textbf{\textit{Оформление курсовой работы:}} \\[0.5cm]
        Расчетно-пояснительная записка на \underline{\hspace{2cm}} листах формата А4. \\[0.5cm]
        \underline{\hspace{15cm}} \\
        \underline{\hspace{15cm}} \\
        \underline{\hspace{15cm}} \\
        \underline{\hspace{15cm}} \\[1.5cm]
        Дата выдачи задания «\und{13}» \und{сентября} \und{2024} г. \\[2cm]
        \textbf{Руководитель курсовой работы} \hfill \underline{\hspace{5cm}} \hfill \und{В.А. Федулов} \\
        \hspace{8cm} (Подпись, дата) \hspace{2cm} (И.О.Фамилия) \\[1.5cm]
        \textbf{Студент} \hfill \underline{\hspace{5cm}} \hfill \und{М.В. Ерофеев} \\
        \hspace{8cm} (Подпись, дата) \hspace{2cm} (И.О.Фамилия)
    \end{flushleft}

\end{titlepage}
\tableofcontents

\newpage
\section*{Введение}
Данная курсовая работа посвящена нахождению оптимальных параметров артиллерийского орудия путем решения обратной задачи внутреннй баллистики идите нахуй всееееееееееееееееее

\end{document}